\section{Electrochemistry}

\subsection{Oxidation States and Redox}

\subsubsection{Oxidation state rules}
\[
\begin{array}{l l}
\toprule
\text{free element} & 0\\
\text{monoatomic ion} & \text{ion charge}\\
\text{oxygen in most compounds} & -II\\
\text{hydrogen in most compounds} & +I\\
\text{halogens in binary salts} & -I\\
\text{sum in molecule or ion} & \text{overall charge}\\
\bottomrule
\end{array}
\]
Oxidation states are bookkeeping numbers used to identify oxidation and reduction.

\subsubsection{Oxidation and reduction}
\[
\text{oxidation: loss of electrons}, \qquad
\text{reduction: gain of electrons}
\]
Oxidation state increases during oxidation and decreases during reduction.

\subsubsection{Basic-solution redox balancing}
\[
\ce{3SO3^2- + 2MnO4- + H2O -> 3SO4^2- + 2MnO2 + 2OH-}
\]
This lecture example conserves atoms and charge in basic solution. It is a reusable check for half-reaction balancing.

\subsection{Galvanic Cells}

\subsubsection{Anode and cathode}
\[
\text{anode: oxidation}, \qquad \text{cathode: reduction}
\]
In a galvanic cell, electrons flow from anode to cathode through the external circuit.

\subsubsection{Cell diagram notation}
\[
\ce{Zn(s) | Zn^2+(aq) || Cu^2+(aq) | Cu(s)}
\]
A single vertical line denotes a phase boundary and the double line denotes a salt bridge.

\subsubsection{Standard cell potential}
\[
E^\circ_{\mathrm{cell}}=E^\circ_{\mathrm{cathode}}-E^\circ_{\mathrm{anode}}
=E^\circ_{\mathrm{red}}(\text{reduction})-E^\circ_{\mathrm{red}}(\text{oxidation})
\]
Standard reduction potentials are used as tabulated. The more positive reduction potential is reduced in a spontaneous galvanic cell.

\subsubsection{Daniell cell potential}
\[
\ce{Zn^2+ + 2e- -> Zn}\quad E^\circ=-0.76\,\mathrm V
\]
\[
\ce{Cu^2+ + 2e- -> Cu}\quad E^\circ=+0.34\,\mathrm V
\]
\[
E^\circ_{\mathrm{cell}}=0.34-(-0.76)=1.10\,\mathrm V
\]
Zinc is oxidized and copper(II) is reduced in the Daniell cell.

\subsection{Free Energy and Equilibrium}

\subsubsection{Electrical work and Gibbs free energy}
\[
\Delta G^\circ=-nFE^\circ_{\mathrm{cell}}
\]
\(n\) is moles of electrons transferred and \(F=9.647\times 10^4\,\mathrm{C\,mol^{-1}}\). Positive \(E^\circ_{\mathrm{cell}}\) gives negative \(\Delta G^\circ\).

\subsubsection{Cell potential and equilibrium constant}
\[
nFE^\circ_{\mathrm{cell}}=RT\ln K
\]
At \(298\,\mathrm K\), a larger standard cell potential corresponds to a larger equilibrium constant.

\subsection{Nernst Equation}

\subsubsection{Nernst equation}
\[
E_{\mathrm{cell}}=E^\circ_{\mathrm{cell}}-\frac{RT}{nF}\ln Q
\]
The cell potential changes with reaction quotient \(Q\). Products in \(Q\) lower the potential when they increase.

\subsubsection{Nernst equation at 298 K}
\[
E_{\mathrm{cell}}=E^\circ_{\mathrm{cell}}-\frac{0.0592\,\mathrm V}{n}\log Q
\]
This base-10 form is used for room-temperature calculations.

\subsubsection{Concentration cell}
\[
E_{\mathrm{cell}}=\frac{0.0592\,\mathrm V}{n}\log\left(\frac{c_{\mathrm{high}}}{c_{\mathrm{low}}}\right)
\]
% Added from exam coverage
Identical electrodes give \(E^\circ_{\mathrm{cell}}=0\). The spontaneous direction reduces ions in the more concentrated half-cell and dissolves metal in the dilute half-cell.

\subsubsection{Solubility product from cell potential}
\[
\Delta G^\circ=-nFE^\circ_{\mathrm{cell}}=-RT\ln K
\]
Electrochemical measurements can determine equilibrium constants, including \(K_{sp}\), when the cell reaction is written for the dissolution or precipitation equilibrium.
